\section{Conclusion}
We revised GNN-AdaptGraph into a cost-explicit adaptive runtime for partition-parallel
processing of dynamic graphs and evaluated it on real execution with established baselines.
Thread reallocation, implemented for real, is the most effective adaptation for
activity-driven imbalance. ACGR, with a communication-volume objective and an explicit balance
constraint, repairs degraded partitions far better than label propagation on sparse graphs.
A GraphSAGE controller trained on outcome-based labels makes better per-decision choices than
a tuned rule. End to end, however, it performs on par with simple heuristics. In particular,
it does not learn when global repartitioning pays off under steady drift, where periodic
METIS remains best. The honest takeaway is a hybrid: a cheap thread-adjustment trigger
combined with drift-triggered global repartitioning. Learning the latter reliably is
the open problem.

Future work will address this gap with longer-horizon objectives (e.g.\ reinforcement
learning over whole streams), training on recorded traces of the target workload, and features
that estimate the benefit of repartitioning cheaply. We will also evaluate on a real
distributed cluster with larger graphs.
